\section*{الفصل \ref{ch:b2:blood-circulation} --- الدم والجملة الدورانية}

\begin{solution}{exo:b2:blood-circulation:1}
\omterm{def:b2:blood-circulation:blood}{البلازما} \qty{55}{\%} (ماء، و\qty{70}{g/L} من البروتين، وأملاح،
ومغذيات، وغازات، وهرمونات)؛ والخلايا \qty{45}{\%}. والكريات الحمر $5\times
10^{12}$ في اللتر، وهي أقراص \qty{7}{\micro m}، لنقل الأكسجين؛ والكريات البيض
$7\times 10^{9}$ في اللتر، بقياس \qtyrange{10}{15}{\micro m}، للمناعة؛ والصفيحات
$3\times 10^{11}$ في اللتر، وهي شظايا \qty{2}{\micro m}، للإرقاء.
\end{solution}

\begin{solution}{exo:b2:blood-circulation:2}
السمكة: قلب $\to$ غلاصم $\to$ جسم $\to$ قلب، أي دارة واحدة، ويتلقى
الجسم \omterm{def:b2:blood-circulation:blood}{الدم} بالضغط المنخفض الباقي بعد الغلاصم. والثديي:
قلب أيمن $\to$ رئتان (\qty{25}{mmHg}) $\to$ قلب أيسر $\to$ جسم
(\qty{120}{mmHg}) $\to$ قلب أيمن. والدارة المضاعفة تعطي
الجسم ضغطًا عاليًا دون تعريض الرئتين له، وتتيح تنظيم
الدارتين على حدة.
\end{solution}

\begin{solution}{exo:b2:blood-circulation:3}
\omterm{def:b2:blood-circulation:vessels}{الشريان}: جدار سميك بطبقات مرنة وعضل --- يوصّل \omterm{def:b2:blood-circulation:blood}{الدم}
بضغط عالٍ وينعّم النبضة. \omterm{def:b2:blood-circulation:vessels}{والشريّن}: جداره في معظمه
عضل أملس --- يضبط المقاومة ويوزع الجريان. \omterm{def:b2:blood-circulation:vessels}{والشعيرة}:
بثخن خلية بطانية واحدة --- للتبادل. \omterm{def:b2:blood-circulation:vessels}{والوريد}: جدار رقيق قابل
للتمدد بدسامات --- يعيد \omterm{def:b2:blood-circulation:blood}{الدم} بضغط منخفض ويخزن معظمه.
\end{solution}

\begin{solution}{exo:b2:blood-circulation:4}
$70\times 72 = \qty{5}{L/min}$، أي \qty{7200}{L} في اليوم، في مقابل
\qty{5}{L} من \omterm{def:b2:blood-circulation:blood}{الدم}: فيمر الحجم عبر القلب \num{1440}
مرة في اليوم. ولا عضو يستطيع صنع سبعة أطنان من \omterm{def:b2:blood-circulation:blood}{الدم} أو استهلاكها في
اليوم؛ فلا بد أن يعود \omterm{def:b2:blood-circulation:blood}{الدم} نفسه --- أي إنه يدور.
\end{solution}

\begin{solution}{exo:b2:blood-circulation:5}
$\Delta P = 8\eta LQ/\pi r^{4} = 8\times 3\times 10^{-3}\times
0.4\times 8.3\times 10^{-5}/(\pi\times 2.44\times 10^{-8}) =
\qty{10}{Pa}$، أي \qty{0.08}{mmHg}: فالأبهر لا يكلف شيئًا؛ و
الضغط يُنفق في \omterm{def:b2:blood-circulation:vessels}{الشرينات}.
\end{solution}

\begin{solution}{exo:b2:blood-circulation:6}
$(15/12)^{4} = 2.4$: فالمقاومة ترتفع 2.4 ضعفًا، والجريان يهبط إلى
\qty{41}{\%}. و$(15/20)^{4} = 0.32$: فالمقاومة تهبط إلى الثلث، والجريان
يرتفع 3.2 ضعفًا.
\end{solution}

\begin{solution}{exo:b2:blood-circulation:7}
الأبهر $83/3 = \qty{28}{cm/s}$؛ والشعيرات $83/3000 =
\qty{0.028}{cm/s}$؛ والعبور $0.1/0.028 = \qty{3.6}{s}$. وفي الجهد
يرتفع النتاج خمسة أضعاف ومساحة الشعيرات المفتوحة ثلاثة
أضعاف ربما، فيهبط زمن العبور إلى نحو \qty{2}{s} --- وهو لا يزال
كافيًا لمغادرة الأكسجين.
\end{solution}

\begin{solution}{exo:b2:blood-circulation:8}
الطرف الشرياني $35 - 25 = \qty{+10}{mmHg}$ (رشح)؛ والطرف الوريدي $15
- 25 = \qty{-10}{mmHg}$ (إعادة امتصاص). ومع $P_c = 28$ عند الطرف
الوريدي يصير الصافي $+3$: فيرشح السائل على الطول كله ولا يُعاد
امتصاص شيء --- أي وذمة، وهي كاحلا قصور القلب المنتفخان.
\end{solution}

\begin{solution}{exo:b2:blood-circulation:9}
$40/66000 = \qty{0.61}{mol/m^3}$؛ و$\pi = 0.61\times 8.314\times 310 =
\qty{1570}{Pa} = \qty{11.8}{mmHg}$؛ ومع نصف الألبومين، \qty{5.9}{mmHg}.
والصوديوم يعبر جدار \omterm{def:b2:blood-circulation:vessels}{الشعيرة} بحرية، فيكون تركيزه
نفسه على الجانبين ولا يمارس ضغطًا تناضحيًا عبره.
\end{solution}

\begin{solution}{exo:b2:blood-circulation:10}
$RC = \qty{2}{s}$: ومنه $120e^{-0.4} = \qty{80}{mmHg}$. ومع $C = 1$: $RC = 1$،
و$120e^{-0.8} = \qty{54}{mmHg}$ --- وحجم الضربة نفسه يرفع
الضغط الانقباضي أكثر في أبهر قاسٍ. \omterm{thm:b2:blood-circulation:windkessel}{وضغط النبض} الواسع
يعني ذروة أعلى على القلب أن يقذف ضدها، وشغلًا أكثر
وأكسجينًا أكثر للنتاج نفسه.
\end{solution}

\begin{solution}{exo:b2:blood-circulation:11}
في الراحة: $90/83 = \qty{1.08}{mmHg.s/mL}$؛ وفي الجهد: $105/417 =
\qty{0.25}{mmHg.s/mL}$، أي هبوط بأربعة أضعاف. ولمّا كان $R \propto r^{-4}$،
ارتفع $r$ بمقدار $4^{1/4} = 1.41$: أي إن \omterm{def:b2:blood-circulation:vessels}{الشرينات} اتسعت
\qty{41}{\%} وسطيًا.
\end{solution}

\begin{solution}{exo:b2:blood-circulation:12}
الاستمرارية تثبّت السرعة عند كل مستوى بالمقطع
الكلي، والشجرة مبنية بحيث يكون المقطع ألف
ضعف مقطع الأبهر في الشعيرات وصغيرًا من جديد في
الأوردة؛ \omterm{thm:b2:blood-circulation:poiseuille}{وقانون بوازوي} يضع المقاومة في \omterm{def:b2:blood-circulation:vessels}{الشرينات}،
حيث يستطيع العضل تغييرها، ويترك الأوعية الواسعة شبه خالية
من الفقد. أما \omterm{def:b2:blood-circulation:circuits}{الدوران المفتوح} فلا شعيرات له تبطئ \omterm{def:b2:blood-circulation:blood}{الدم} في
موضع محدد، ولا أوعية يضبط فيها مقاومة، ولا سبيل
لتوجيه الجريان إلى عضو واحد --- فلا يستطيع إلا التقليب.
\end{solution}

\begin{solution}{pb:b2:blood-circulation:1}
\textbf{1.} $70\times 72 = \qty{5.0}{L/min}$؛ أي \qty{7300}{L} في اليوم.
\textbf{2.} $7300/5 \approx 1450$ مرة.
\textbf{3.} ربع \qty{57}{g} عند 72 نبضة في الدقيقة: أي \qty{14}{g}
$\times 4320 = \qty{62}{kg}$ في الساعة --- أي نحو وزن رجل.
\textbf{4.} لا طعام يستطيع أن يوفّر، ولا نسيج أن يستهلك، وزن رجل
من \omterm{def:b2:blood-circulation:blood}{الدم} كل ساعة؛ والمخرج الوحيد أن يعود \omterm{def:b2:blood-circulation:blood}{الدم} نفسه
إلى القلب.
\textbf{5.} خيط يُربط حول الذراع ينفخ الأوردة أسفله لا
أعلاه، \omterm{def:b2:blood-circulation:blood}{فالدم} في الأوردة يتحرك نحو القلب؛ وبدفع
\omterm{def:b2:blood-circulation:blood}{الدم} في \omterm{def:b2:blood-circulation:vessels}{وريد} نحو اليد، يقف عند الدسامات ولا يمكن
دفعه رجوعًا --- فالدسامات لا تسمح بالجريان إلا نحو القلب.
\textbf{6.} الشعيرات الواصلة بين الشرايين والأوردة؛ ورآها مالبيغي
في رئة الضفدع سنة 1661.
\textbf{7.} المساحة $\pi\times 1.25^{2} = \qty{4.9}{cm^2}$؛ ومنه $83/4.9 =
\qty{17}{cm/s}$.
\textbf{8.} $8\times 3\times 10^{-3}\times 0.4\times 8.3\times
10^{-5}/(\pi\times 2.44\times 10^{-8}) = \qty{10}{Pa}$، أي \qty{0.08}{mmHg}.
\textbf{9.} $R_{1} = 8\times 3\times 10^{-3}\times 10^{-3}/(\pi\times
5.06\times 10^{-20}) = \qty{1.5e14}{Pa.s/m^3}$؛ وعلى التوازي،
$1.5\times 10^{14}/3\times 10^{6} = \qty{5.0e7}{Pa.s/m^3}$.
\textbf{10.} $8.3\times 10^{-5}\times 5.0\times 10^{7} = \qty{4200}{Pa}
= \qty{31}{mmHg}$.
\textbf{11.} الشعيرات المفتوحة $10^{10}$، ومساحة كل منها $\pi(3\times 10^{-6})^{2}
= \qty{2.8e-11}{m^2}$: أي \qty{0.28}{m^2} $= \qty{2800}{cm^2}$؛ ومنه $v =
83/2800 = \qty{0.03}{cm/s}$؛ والعبور $0.1/0.03 = \qty{3.3}{s}$.
\textbf{12.} المقاومة $\times(1/0.9)^{4} = 1.52$: أي \qty{47}{mmHg}؛ وعلى
القلب أن يرفع الضغط الشرياني \qty{16}{mmHg} وإلا هبط
النتاج بمقدار الثلث.
\textbf{13.} $+10$ و$-10$ و$0$ ملم زئبق.
\textbf{14.} الرشح $10\times 2 = \qty{20}{L}$ في اليوم؛ وإعادة الامتصاص
$8.5\times 2 = \qty{17}{L}$.
\textbf{15.} \qty{3}{L} من اللمف في اليوم.
\textbf{16.} $\pi_c = \qty{12.5}{mmHg}$: فالصافي $+22.5$ عند الطرف
الشرياني و$+2.5$ عند الطرف الوريدي --- أي رشح في كل مكان، ولا
إعادة امتصاص: أي وذمة.
\textbf{17.} الصافي $+10$ و$+3$، ومتوسطه $+6.5$: أي رشح على
امتداد \omterm{def:b2:blood-circulation:vessels}{الشعيرة} كلها، بنحو \qty{26}{L} في اليوم؛ ولا تستطيع اللمفيات حمله
فيتراكم السائل في الأنسجة، بدءًا بالقدمين.
\textbf{18.} $4\times 10^{10}\times 2\pi\times 3\times 10^{-6}\times
10^{-3} = \qty{750}{m^2}$؛ و$t = x^{2}/2D = (2\times 10^{-5})^{2}/
(2\times 10^{-9}) = \qty{0.2}{s}$.
\textbf{19.} $(93 - 3)/83 = \qty{1.08}{mmHg.s/mL}$.
\textbf{20.} $RC = \qty{2.2}{s}$؛ ومنه $120e^{-0.8/2.2} = 120\times 0.69 =
\qty{83}{mmHg}$.
\textbf{21.} $RC = \qty{1.1}{s}$؛ ومنه $120e^{-0.73} = \qty{58}{mmHg}$: فيتسع
\omterm{thm:b2:blood-circulation:windkessel}{ضغط النبض} من 37 إلى \qty{62}{mmHg}.
\textbf{22.} $120e^{-0.3/2.2} = \qty{105}{mmHg}$: فعند معدل سريع لا
يكاد الضغط يهبط بين النبضات.
\textbf{23.} $C\Delta P = 2\times 20 = \qty{40}{mL}$ مخزونة؛ أما
\qty{30}{mL} الأخرى فتنفد عبر \omterm{def:b2:blood-circulation:vessels}{الشرينات} خلال الانقباض نفسه.
\textbf{24.} البطين المنقبض يعصر أوعيته فيغلقها
في الانقباض، فتُروى عضلة القلب في الانبساط، بالضغط
الذي يصونه الأبهر المرتد؛ والأبهر القاسي الذي يدع
الضغط الانبساطي ينهار يجوّع القلب بين النبضات.
\textbf{25.} النتاج \qty{5}{L/min}؛ والهبوط الشرييني \qty{31}{mmHg}؛
والرشح الصافي \qty{3}{L} في اليوم إلى اللمف؛ والانبساطي \qty{83}{mmHg}
من أجل $C = 2$ و\qty{58}{mmHg} من أجل $C = 1$.
\end{solution}
